\section{Fording, and what an untyped target changes}
\label{sec:fording}

This section reviews the type-theoretic situation the paper starts from,
fixes the reading of \emph{fording} used throughout, and identifies which
parts of the classical machinery an untyped first-order target renders
unnecessary --- and which soundness obligations it creates instead.  It
contains no formal development; the definitions begin in
\cref{sec:calculus}.

\subsection{Indexed families and their eliminations}
\label{sec:fording-families}

An inductive family in the sense of Dybjer \cite{Dybjer1994} is an
inductive definition $I : \tup{\sigma} \to s$ whose constructors target
\emph{instances} $I\,\tup{t}(\tup{a})$ determined by the constructor's own
arguments.  The paradigmatic informative examples are
\[
  \csvnil : \vect\,A\ \csO, \qquad
  \csvcons : \forall n{:}\nat.\ A \to \vect\,A\ n \to \vect\,A\ (\csS\,n),
\]
and the paradigmatic propositional ones are the equality family and
$\mathsf{le}$.  Three features distinguish families from plain datatypes
in a way that any translation must confront.

\begin{enumerate}[leftmargin=2em]
\item \emph{Constructor typing constrains indices.}  A constructor
  inhabits only the instances of the form $I\,\tup{t}(\tup{a})$; an
  arbitrary instance $I\,\tup{u}$ is inhabited by $\wcns{c}$-forms only
  when the constraint $\tup{u} = \tup{t}(\tup{a})$ is solvable.
\item \emph{Case analysis specializes indices.}  A dependent match on a
  scrutinee of type $I\,\tup{u}$ must reconcile $\tup{u}$ with each
  constructor's $\tup{t}_i(\tup{y}_i)$.  Type-theoretic elaborations do so
  by first-order unification with constructors as the rigid symbols
  \cite{GoguenEtAl2006}; branches whose constraints are refutable are
  discharged, and solvable constraints instantiate the branch.
\item \emph{Propositional families may or may not eliminate into
  computations.}  CIC permits a match on a $\Prop$-sorted scrutinee in an
  informative position only for empty and \emph{singleton} inductives
  (one constructor, all arguments propositional, parameters aside)
  \cite{Rocq,Letouzey2003}; the equality family qualifies because its
  constructor has no arguments beyond parameters.  The criterion is known
  to be simultaneously too strong and too weak
  \cite{GilbertEtAl2019}: it rejects families whose non-parameter
  arguments are \emph{forced} by the indices, and accepting it together
  with proof irrelevance makes uniqueness of identity proofs (UIP) hold.
\end{enumerate}

\subsection{Fording}
\label{sec:fording-fording}

McBride's thesis \cite[\S 3.5]{McBride1999} introduces the constraint
technique this paper is built on: to use an elimination rule whose motive
ranges over the whole ``aperture'' $I\,\tup{\imath}$ at a more specific
instance, one abstracts the full aperture and \emph{constrains} it with
pointwise equations
$\tup{\imath} \simeq \tup{u}$ --- ``the patterns to which it applies are
subject to equational constraints which recover their specificity.''
Goguen, McBride and McKinna's elimination of dependent pattern matching
\cite{GoguenEtAl2006} runs on the same device: in their specialization by
unification, the method for constructor $\wcns{c}$ receives hypotheses
$\tup{u} \simeq \tup{t}_c(\tup{y})$ and $\wcns{c}\,\tup{y} \simeq x$ ---
``the equations on the indices are exactly those we must unify to allow
the instantiation.''  At the level of declarations the transformation is
folklore, stated compactly by Zhang \cite{Zhang2023}: a constructor
$\wcns{c} : X \to F\,a$ fords to
$\wcns{c}^{\mathrm{Ford}} : \forall a'.\ (a = a') \to X \to F\,a'$, after
which ``$F$ is no longer indexed'' --- the family has become a
\emph{parameterized} type whose constructors carry equational premises,
and index specialization has become ordinary propositional reasoning that
a consumer may perform whenever convenient.

This paper uses fording in three places, always as a lens and never as a
literal re-declaration:

\begin{itemize}[leftmargin=2em]
\item in the \emph{source calculus}: the branches of a case analysis on
  $I(\tup{u})$ bind proofs of the index equations
  $\tup{u} = \tup{t}_i(\tup{y}_i)$ (\cref{fig:terms}), which is
  what lets dependently typed idioms --- impossible branches, results
  depending on indices --- be written with first-order motives;
\item in the \emph{emitted theory}: the inversion axiom of a family
  quantifies over an arbitrary index tuple and disjoins, per constructor,
  the payload \emph{and the index equations}
  (\cref{def:theory}) --- the first-order image of the forded
  declaration;
\item in the \emph{semantics}: the interpretation of a family is a set of
  (index tuple, element) pairs, and the index-reflection lemma
  (\cref{lem:index-reflection}) states that membership holds exactly when
  a constructor decomposition with matching index denotations exists ---
  fording read as a property of one model.
\end{itemize}

One point deserves emphasis, because it explains why no
``forded re-declaration'' appears anywhere in the formal development.
Read as a CIC declaration, the forded family has its erstwhile indices as
\emph{non-uniform parameters}: recursive constructor arguments still
occur at other instantiations ($\csvcons$ at $n$, conclusion at
$\csS\,n$).  Parameter-versus-index is a property of constructor
\emph{conclusions} only, and elimination confronts only conclusions:
since every forded constructor returns the generic instance, case
analysis needs no index unification, and each branch simply receives its
equations as hypotheses.  Constraint discharge is pay-as-you-go, one
constructor layer per match.  In the untyped first-order output the
parameter/index distinction vanishes entirely --- there are only function
symbols, predicates and equations --- which is why fording can remain a
per-occurrence, per-constructor lens.

\subsection{What the untyped target dissolves}
\label{sec:fording-dissolves}

The type-theoretic literature attaches three warning labels to fording
and to dependent case analysis generally.  All three are analysed
precisely in the body of the paper; we preview the outcome.

\paragraph{Heterogeneous equality.}  Telescopic index equations relate
terms whose types agree only after earlier equations are substituted:
$v_1 = v_2$ at types $\vect\,A\,n_1$ and $\vect\,A\,n_2$ cannot even be
\emph{stated} homogeneously before $n_1 = n_2$ is substituted.  McBride's
heterogeneous $\simeq$ exists to state such equations, and its eliminator
is exactly K \cite{McBride1999,CockxEtAl2014}.  In an untyped target the
formation problem does not exist: equality is a single binary predicate
on one universe, and telescopic equations are stated pointwise
(\cref{def:theory}).  The paper's source calculus sidesteps the issue
differently --- its index telescopes are non-dependent
(\cref{def:env}), so source-level equations are homogeneous --- and
\cref{sec:limitations} discusses the residual gap.

\paragraph{K, UIP, and the singleton criterion.}  The without-K analysis
of pattern matching \cite{CockxEtAl2014,CockxDevriese2018} isolates
exactly two unification steps that require K: deletion of reflexive
equations, and injectivity at indices that are non-linear or not
self-unifiable; Gilbert et al.\ \cite{GilbertEtAl2019} show that
eliminating an indexed propositional singleton under proof irrelevance
\emph{is} UIP.  Every one of these obstructions is a proof-relevance
obstruction.  The model of \cref{sec:semantics} is proof-irrelevant by
construction, hence validates UIP outright, so the \emph{unrestricted}
rule set is sound for the translation: deletion, non-linear injectivity
and singleton collapse need no side conditions
(\cref{sec:formulas-uip}).  What survives of the without-K literature is
its \emph{negative} space: discrimination only for fully applied
constructor forms (under-applied conflict is refutable under functional
extensionality --- Counterexample~35 of \cite{CockxDevriese2018}), and no
injectivity for type formers (\cref{rem:hur}).

\paragraph{Green slime.}  Defined functions in index positions make
unification-based elaboration stuck \cite{McBride2014}: $0 \simeq
\mathsf{plus}\,m\,n$ is neither a clash nor a solution.  The standard
type-theoretic workarounds (ford the offending index, relational
re-expression, derived eliminators) all amount to \emph{keeping the
equation as data}.  The translation does exactly that, uniformly, for
every index: the emitted equation $z \feq
\Ihat{\mathsf{plus}}(m,n)$ sits in an inversion disjunct or a guard
formula next to the defining equations of $\Ihat{\mathsf{plus}}$, and a
saturation prover performs specialization by unification as ordinary
equational reasoning --- including closing branches that elaboration-time
unification cannot (a refutation of $0 \feq \mathsf{plus}(m,n)$ under
$m \feq \csS(k)$ is one paramodulation step away from the defining
equations).  Nothing in the soundness proof distinguishes rigid from
non-rigid index terms; decidability is neither claimed nor needed
(\cref{rem:slime}).

\subsection{What the untyped target does not dissolve}
\label{sec:fording-obligations}

Three families of documented accidents mark the soundness obligations
that an untyped target \emph{creates}; the formal development treats each
as a theorem rather than a guideline.

First, \emph{exhaustiveness}: in a one-universe target, ``every $x$ is a
constructor form'' bounds the cardinality of the whole universe
\cite[\S 2.3]{MengPaulson2008}, and the monotonicity analysis of
\cite{BlanchetteEtAl2016} shows such sentences are exactly the ones whose
guards can never be erased.  For indexed families the situation is
strictly worse: for every family that specialises an index to a
constructor form --- vectors, finite sets, and their relatives --- an
\emph{unguarded} inversion axiom is not merely false in the intended
model but inconsistent with the discrimination axioms the translation
must also emit, its own index equations identifying distinct
constructor forms (\cref{prop:unguarded-inconsistent});
enumeration-shaped families detonate by cardinality instead
(\cref{rem:unguarded-general}).

Second, \emph{index erasure in atoms}: a type-indexed notion encoded by
an index-blind target notion lets facts established at one index leak to
another.  The \fstar{} unsoundness \#1542 \cite{FStar1542} is the
documented instance --- extensional function equality at domain
$\mathsf{nat}$, encoded into the SMT solver's untyped $=$, leaked to
domain $\mathsf{int}$ and proved $\mathsf{False}$.  The analogue inside
this paper's design space is a guard predicate that forgets its index
arguments; \cref{prop:index-dropping} shows the resulting translation
proves a semantically false goal.  The translation of
\cref{sec:translation} therefore keeps the indices as arguments of every
guard atom, and its native untyped $\feq$ is sound for a different reason
than \fstar{}'s was unsound: the \emph{semantics of source equality} in
the model is itself untyped conversion-class equality, so no typed
reading exists to be violated --- at the documented price that the model
is intensional and functional extensionality is not validated
(\cref{rem:funext-leak}).

Third, \emph{degenerate instances}: axioms asserting existence or
structure at empty instances --- selector totality, box/unbox roundtrips,
rank orderings --- have produced inconsistencies in deployed systems
(Verus \#1366 \cite{Verus1366}).  The translation emits no rank axioms,
no selectors and no unboxing roundtrips; at an empty instance the guard
atom is simply false in the model and every structural axiom is
vacuously satisfied there, while refutations of inhabitation
(``$\Fin(0)$ is empty'') become one-step consequences of the guarded
inversion axiom and discrimination (\cref{ex:fin-empty}).
